\section{Positive-time covariance and macroscopic entropy inputs}
\label{prof:positive-inputs}

We establish the fixed-positive-time estimates used to start the backward
induction, and the estimates for changes of mean needed at macroscopic
entropy.  The matrix in a heat-bath conditional mean is
$J^\circ=J-\operatorname{diag}(J_{11},\ldots,J_{nn})$.  Thus, in this
section, put
\[
 \begin{gathered}
 m(y)=\tanh y,\quad q(y)=n^{-1}\|m(y)\|^2,\quad b(y)=1-q(y),\\
 V_y=\operatorname{diag}(1-m_i(y)^2),\qquad
 F_h(y)=y-h+(b(y)I-J^\circ)m(y).
 \end{gathered}
\]
For a nonnegative diagonal matrix $A$ define
\[
 \mathcal H(y,A)=I+A^{1/2}
 \left(b(y)I-J^\circ-\frac{2m(y)m(y)^{\mathsf T}}n\right)A^{1/2}.
\]
The relevant stability property, with fixed positive constants, is
\begin{equation}\label{prof:eq26}
 \left.
 \begin{gathered}
 \|F_h(y)\|\le\rho_0\sqrt n,\qquad
 0\le A\le(1+\eta_0)I,\\
 n^{-1/2}\|A-V_y\|_{\mathrm F}\le\epsilon_A
 \end{gathered}\right\}
 \quad\Longrightarrow\quad \mathcal H(y,A)\succeq c_0 I.
\end{equation}
Only diagonal $A$ are quantified in this implication.

The companion high-temperature article
\cite{AcceptedGap} proves the requisite deterministic residual
identities and posterior consequences.  Its probabilistic matrix and
stable-field statements are stated for $j=\beta^2<1$; they do not directly
supply \eqref{prof:eq26} at $j=1$.  We give the endpoint replacements below.
In particular, none of the following arguments uses its spectral-gap
theorem at the critical temperature.

\subsection{Inverse words with a strict norm--trace margin}

For a diagonal $A\ge0$ write $\chi_A=n^{-1}\operatorname{Tr}A$ and set
\[
 S_A(z;M)=I+z^2\chi_AA-zA^{1/2}MA^{1/2},\qquad
 K_A(z;M)=(I-zAM+z^2\chi_AA)^{-1}.
\]
Fix $a_+<\infty$ and $\kappa\in(0,1)$, and define the closed parameter set
\[
 \mathcal A_\kappa=
 \{A\text{ diagonal}:0\le A\le a_+I,\quad \|A\|\chi_A\le1-\kappa\}.
\]
An ordinary word is a finite product of bounded diagonal matrices and
copies of $M$.  An inverse word has, in addition, at most one factor
$K_A(1;M)$.  Its diagonal prediction is obtained by noncrossing
contractions $MDM\mapsto n^{-1}\operatorname{Tr}(D)I$.  For an inverse
factor use its formal series in $z$ before contracting; the prediction
at $z=1$ is well defined by the polynomial assertion in the next lemma.
The associated normalized trace prediction is cyclic: rotating the
factors of a word does not change
$n^{-1}\operatorname{Tr}\mathsf P(Q)$, coefficient by coefficient
in the formal variable $z$.

\begin{lemma}[Restricted endpoint word estimates]
\label{prof:endpoint-words}
Let $W$ be GOE with off-diagonal variance $1/n$ and diagonal variance
$2/n$, and let $W^\circ$ be its zero-diagonal version.  Fix a word length
$L$, a bound $M_0$ on ordinary diagonal factors, and an inverse margin
$c>0$.  Except on an event of probability $Ce^{-c_1n}$, simultaneously
for $M=W,W^\circ$ and all such words $Q$,
\[
 \|Q\|+
 \left(\sum_{i\ne j}|Q_{ij}|^4\right)^{1/4}
 +\|\operatorname{diag}(Q-\mathsf P(Q))\|_2\le C.
\]
For inverse words this assertion is restricted to $A\in\mathcal A_\kappa$
and $S_A(1;M)\succeq cI$. All diagonal parameters and permitted
word choices within the fixed bounds may be selected after observing
$M$; the bounds $L,M_0,c,a_+,\kappa$ are fixed beforehand.
Constants depend on the displayed fixed bounds, including $\kappa$.
The prediction for a word with $r$ ordinary matrix letters outside its
inverse is a polynomial in $z$ of degree at most $r$.

For each deterministic $A\in\mathcal A_\kappa$ there is a constant
$c_\kappa>0$ such that
$S_A(z;W)\succeq c_\kappa I$ for all $z\in[0,1]$ except with probability
$Ce^{-c_2n}$.  For any fixed finite family of deterministic bounded vectors,
the bilinear deterministic equivalents are, to any fixed positive
accuracy,
\[
 G_A\simeq A,\qquad G_AW\simeq\chi_A A,\qquad
 WG_AW\simeq\chi_A I+\chi_A^2 A,
 \quad G_A=A^{1/2}S_A(1;W)^{-1}A^{1/2},
\]
with the same exponential form of exceptional probability.  Here each
comparison concerns only the specified bilinear forms.
\end{lemma}

\begin{proof}
We give the endpoint changes to Lemmas 3.1--3.4 of
\cite{AcceptedGap}, including the uniformity needed below. The proof
separates fixed parameters from adaptive ones. A strict norm--trace bound
gives the deterministic-diagonal path margin and controls expectations.
Chaining makes the word errors simultaneous over the diagonal parameters.
An inverse chosen after observing the matrix is used only where its own
stated positive margin holds.
For $\|p\|=1$ and $w=A^{1/2}p$, Gaussian comparison bounds the expected
supremum of
$zw^{\mathsf T}Ww-z^2\chi_A\|w\|^2$ by that of
\[
 2zT_g v-z^2\chi_Av^2,
 \qquad 0\le v\le\|A\|^{1/2},\qquad
 T_g=n^{-1/2}\|A^{1/2}g\|.
\]
Indeed the increment difference between the comparison process
$2n^{-1/2}\|w\|g^{\mathsf T}w$ and $w^{\mathsf T}Ww$ is
\[
 \frac2n\left[(\|w\|^2-\|w'\|^2)^2
       +2(\|w\|\|w'\|-w^{\mathsf T}w')^2\right]\ge0.
\]
Also $\mathbb E T_g^2=\chi_A$ and
$\operatorname{Var}(T_g^2)\le2a_+^2/n$, so replacing $T_g$ by
$\sqrt{\chi_A}$ costs $o(1)$ uniformly in $A$.

The supremum then has
integrand $2d-d^2$, where
\[
 0\le d=z\sqrt{\chi_A}v\le\sqrt{\|A\|\chi_A}
 \le\sqrt{1-\kappa}.
\]
Its distance below one is at least
$\delta_\kappa=(1-\sqrt{1-\kappa})^2$.
The supremum is $a_+$-Lipschitz in $W$ for the Frobenius metric.
Gaussian concentration, the GOE norm tail, and a fixed mesh in $z$
give the asserted path margin, for example with
$c_\kappa=\delta_\kappa/4$.

To prove expectations without conditioning on a matrix event, project
$W$ onto a fixed operator-norm ball and replace the symmetric inverse
by a smooth bounded functional-calculus function agreeing with $x^{-1}$
on an interval containing all spectra with either required margin.
Use
\[
 K_A=I+A^{1/2}S_A^{-1}A^{1/2}(zW-z^2\chi_AI),\qquad
 K_AA=A^{1/2}S_A^{-1}A^{1/2}.
\]
This defines bounded truncated words.  For fixed length their
Frobenius Lipschitz constants in $W$ are bounded, and the increment
between parameter tuples at sup-norm distance $\Delta$ has
Frobenius Lipschitz constant at most $C\sqrt\Delta$.

These follow by the
product rule, $\|A^{1/2}-B^{1/2}\|\le\|A-B\|^{1/2}$, and the
Fourier formula for differentiating a smooth matrix function.  Thus
entry variances are $O(n^{-1})$, normalized-trace variances are
$O(n^{-2})$, and fourth centered entry moments are $O(n^{-2})$.
These bounds gain respectively $\Delta,\Delta,\Delta^2$ for increments.

The fixed-parameter agreement event has exponentially small complement.
Gaussian integration by parts therefore gives, with errors $O(n^{-1})$
uniformly in $i,z,A$, the same loop equations as in Lemma 3.3 of the
companion article.  Writing $a_i=A_{ii}$,
$k_i=\mathbb E(K_A)_{ii}$, $r=n^{-1}\sum_i a_i k_i$, and $h=r-\chi_A$,
these equations are
\[
 \mathbb E(WK_A)_{ii}=zr k_i+O(n^{-1}),\qquad
 k_i=1+z^2a_i h k_i+O(n^{-1}).
\]
All expectations here may use the bounded truncations.  Their
integration-by-parts error is exponentially small times a fixed
polynomial in $n$, because both true and formal differentials are
bounded and agree on the path event.

If $|h|\le\delta$, the second
equation gives $k_i=1+O(\delta+n^{-1})$.  Consequently
\[
 \left|z^2n^{-1}\sum_i a_i^2k_i\right|
 \le(1+C\delta+C/n)n^{-1}\sum_i a_i^2
 \le(1+C\delta+C/n)(1-\kappa)\le1-\kappa/2
\]
for fixed sufficiently small $\delta$.  Averaging the loop equation
with weights $a_i$ gives $|h|\le C/n$.  Continuity from $h(0)=0$
prevents an exit from $|h|\le\delta$.  It follows that
$k_i=1+O(n^{-1})$ and $\mathbb E(WK_A)_{ii}=z\chi_A+O(n^{-1})$.

For completeness, the higher-word recursion has no new stability
requirement.  Mark an ordinary matrix letter, integrate it by parts,
and let its differentiated partner split the word as $U\boxed W
V\boxed W Q$.  The leading term is
\[
 \mathbb E(UQ)_{ii}\,n^{-1}\mathbb E\operatorname{Tr}V;
\]
the reversed contraction is $n^{-1}\mathbb E(UV^{\mathsf T}Q)_{ii}$
and is $O(n^{-1})$.  Covariance errors are $O(n^{-1})$ by the preceding
variance bounds.  If the derivative meets the inverse, use
$\mathrm dK_A[E]=zK_AAEK_A$.

Splitting then puts one copy of $K_A$
in each smaller word, whose total number of ordinary letters is reduced
by one.  Otherwise it is reduced by two.  The base formal prediction
is $\mathsf P(K_A)=I$: for each nonempty noncrossing-paired chain,
summing over subsets of its originally adjacent pairs gives
$\sum_T(-1)^{|T|}=0$.  Induction proves the polynomial assertion and
entrywise expectation error $O(n^{-1})$.  Every inverse in this
calculation has the original coefficient $A$; no enlarged diagonal
parameter set is used.
For a normalized trace, the noncrossing face formula has no marked
starting position on the trace cycle. Rotating the word permutes its
faces and contractions, proving the stated cyclic invariance
coefficient by coefficient, also after the formal inverse expansion.

To make the result simultaneous, take sup-norm meshes whose points lie
in $\mathcal A_\kappa$ and in the bounded boxes for the other diagonal
parameters.  At scale $2^{-l}$ their cardinalities are at most
$\exp(C_L(l+1)n)$.  The truncated estimates hold in the ambient box,
so mesh increments are legitimate even though no convexity of
$\mathcal A_\kappa$ is assumed.  The diagonal $\ell^2$ and off-diagonal
$\ell^4$ seminorms of centered increments have expectations at most
$C2^{-l/2}$ and concentration scale $C\sqrt{2^{-l}/n}$.

A union bound bounds every increment by
$C\sqrt{l+1}\,2^{-l/2}$, with summable failure probabilities
$Ce^{-c(l+1)n}$.  Sum the increments and the fixed-parameter biases.
Sign conjugation makes all off-diagonal expectations zero.
This proves the asserted simultaneous estimates for $W$.
Since $\|\operatorname{diag}W\|_{\mathrm F}$ is bounded except with
exponentially small probability, the truncated Frobenius Lipschitz
estimate proves the result for $W^\circ$ as well.  Restriction to the
stated positive margin restores the exact inverse factors.

Finally sign conjugation and the loop equations give the diagonal
expectations of $G_A$, $G_AW$ and $WG_AW$. They are respectively $A$,
$\chi_AA$ and $\chi_AI+\chi_A^2A$, to $O(n^{-1})$ in each entry.  Their truncated
bilinear forms are uniformly Lipschitz.  Gaussian concentration proves
the last assertion for each fixed finite list of vectors.
\end{proof}

\begin{lemma}[Restricted determinant bound]\label{prof:endpoint-logdet}
Fix $\kappa,\varepsilon>0$, a rank bound $r_0$, and a norm bound $M_2$.
There is $\gamma_0>0$ such that for each fixed
$0<\gamma\le\gamma_0$ and each fixed $B>0$, with probability at least
$1-e^{-Bn}$ for sufficiently large $n$, simultaneously for
\[
 V\ \text{diagonal},\qquad 0\le V\le I,\qquad
 b=n^{-1}\operatorname{Tr}V\le1-\kappa,
 \qquad \operatorname{rank}E\le r_0,\quad\|E\|\le M_2,
\]
the matrix $L=I+(bI-W)V+E$ satisfies
\[
 n^{-1}\log\det(L^{\mathsf T}L+\gamma I)\le b^2+\varepsilon.
\]
The choices of $\gamma_0$ and the threshold on $n$ may depend on the
fixed margins.  The threshold on $n$ may also depend on $B$.
\end{lemma}

\begin{proof}
For deterministic $V$ use Lemma~\ref{prof:endpoint-words} along
$L_z=I+z^2bV-zWV$.  On its path event,
\[
 \frac{\mathrm d}{\mathrm dz}\frac1n\log\det L_z
 =2zb\,\frac{\operatorname{Tr}(K_V(z)V)}n
       -\frac{\operatorname{Tr}(WK_V(z)V)}n.
\]
The expectation is $zb^2+O(n^{-1})$, by the loop calculation.
Using a fixed norm cutoff in the path event $\mathcal E_V$, integration
from zero to one gives
$\mathbb E[\mathbf1_{\mathcal E_V}2n^{-1}\log\det L_1]
=b^2+O(n^{-1})$.  The inverse norm there is bounded; regularization costs
$C_\kappa\gamma$.  Off this event, the regularized logarithm is bounded
in absolute value by $|\log\gamma|+C+2\log(1+\|W\|)$ and its expected
contribution tends exponentially to zero.  Hence the regularized mean
is at most $b^2+C_\kappa\gamma+o(1)$ uniformly in $V$.

The map $M\mapsto n^{-1}\log\det(M^{\mathsf T}M+\gamma I)$ has
Frobenius Lipschitz constant $1/\sqrt{n\gamma}$.  Gaussian
concentration thus gives a tail $2\exp(-\gamma n^2u^2/4)$ at deviation
$u$.  Choose $\gamma_0$ so $C_\kappa\gamma_0<\varepsilon/8$.
A fixed sup-norm mesh of the restricted diagonal set has size
$\exp(O(n))$.  On $\|W\|\le R$, a mesh spacing depending on
$R,\gamma,\varepsilon$ controls both the logarithm and $b^2$ to
$\varepsilon/4$.  For this last mesh step choose a possibly larger $R$ so that its
exceptional probability is at most $e^{-(B+1)n}$.

This choice does not
change the preceding expectation estimate or its regularizer threshold.  The $n^2$ concentration bound
pays for this mesh for sufficiently large $n$.  Finally a rank-$r_0$
perturbation changes the logarithm by at most
$M_2\sqrt{r_0/(n\gamma)}$.  This proves the statement, including its
arbitrarily large fixed exponential budget.  The exponentially small
path exception was used only in the mean estimate; its exponent does
not limit this final budget.
\end{proof}

The diagonality restriction is essential for this simultaneous
statement. If arbitrary positive contractions chosen after \(W\)
were allowed, take \(V=\mathbf1_{\{W\le0\}}\) and \(E=0\).
Then \(b\to1/2\), so the constraint holds, for example, at
\(\kappa=1/4\), while every eigenvalue of \(L\) on the
range of \(V\) is at least \(1+b\), and the other eigenvalues
equal one. For every \(\gamma>0\),
\[
 \frac1n\log\det(L^TL+\gamma I)
 \ge2b\log(1+b)\longrightarrow\log(3/2)>1/4.
\]
Thus the asserted bound with arbitrarily small \(\varepsilon\)
would fail in that larger class. The application below uses only
the diagonal matrix \(V_y\).

\subsection{Scalar equality and stable fields at positive times}

\begin{lemma}[Endpoint scalar rate]\label{prof:endpoint-scalar}
Let $P$ have finite second moment, put
$q=\mathbb E_P\tanh^2 y$, $b=1-q$, and
$a=\mathbb E_P[(y-d)\tanh y]$.  For $\nu>0$ with $\nu\ge q$,
\[
 D(P\Vert N(d,\nu))\ge\frac{(a-\nu b)^2}{2\nu(\nu+q)}.
\]
Equality holds exactly when $P=N(d,\nu)$.  At every observation time
$s>0$ the consistent equality law, with
$d=s+\mathbb E_P\tanh y$ and $\nu=s+q$, is uniquely
$N(v_*(s),v_*(s))$, where
\[
 v_*(s)=s+\mathbb E_{N(v_*(s),v_*(s))}\tanh y.
\]
On compact positive time intervals this curve is continuous and its
$q$ is bounded away from zero.  The scalar integral and rate conclusions
of Lemma 5.4(i), (ii), and (iv) of \cite{AcceptedGap} therefore
hold at $j=1$ on such intervals.  In particular, outside any fixed
neighborhood of the equality curve, the smoothed scalar integral has
strictly negative rate for sufficiently small smoothing variance.
\end{lemma}

\begin{proof}
Suppose the entropy is finite; then $0<q<1$.  Set
$k=q/\nu\in(0,1]$ and $C=a/\nu$.  The change of variable
$T_u(y)=y-u\tanh y$, $u<1$, gives
\[
 D(P\Vert N(d,\nu))
 \ge u(C-b)-\frac{k u^2}{2}-b[-\log(1-u)-u].
\]
This identity is the entropy change-of-variables calculation in
Lemma 4.1 of the companion article and uses no strict restriction on
$j$.  Put $h(u)=u^2/2-b[-\log(1-u)-u]$.  Since $b=1-q<1$,
\[
 h(u)>0\quad(u\le q,\ u\ne0),\qquad h(q)\ge q^3/6.
\]
For negative $u$ use $-\log(1-u)-u\le u^2/2$; for positive $u\le q$
use the increasing series ratio and
$h(q)=\sum_{l\ge3}q^l/[l(l-1)]$.

If $(C-b)/(1+k)\le q$, substitution proves strict inequality unless
$C=b$.  Otherwise let $C_0=1+kq$ and
$C_1=C_0+\sqrt{(1+k)/3}\,q^{3/2}$.  The choice $u=q$ proves strictness
for $C_0<C<C_1$.  For $C\ge C_1$, variance comparison and
Cauchy--Schwarz give
\[
 2D(P\Vert N(d,\nu))\ge C^2/k-1-\log(C^2/k).
\]
Subtract the target and write
$f(C)=C^2/k-1-\log(C^2/k)-(C-b)^2/(1+k)$.  Direct differentiation gives
\[
 f(C_0)\ge-2q^3/3,\qquad
 f'(C_0)\ge\frac{2q}{1+q},\qquad f''(C)\ge1.
\]
Consequently
\[
 f(C_1)\ge q^3\left[-\frac23+
       \frac{2}{\sqrt3(1+q)\sqrt q}+\frac16\right]>0.
\]
Since $f'$ is positive and increasing after $C_0$, this finishes the
strict cases.  If $C=b$, the target is zero and equality is equivalent
to vanishing entropy.  Gaussian integration by parts verifies equality
for the stated Gaussian law.

For a fixed variance $\nu>0$, the function
$d\mapsto d-\mathbb E_{N(d,\nu)}\tanh y$ is strictly increasing,
onto $\mathbb R$, and zero at $d=0$; its derivative is
$\mathbb E\tanh^2y>0$.  Thus its solution at $s\ge0$ is nonnegative.
Let $H_\nu(d)=\mathbb E_{N(d,\nu)}(\tanh y-\tanh^2y)$.
Gaussian reflection gives $H_\nu(\nu)=0$ and
$\mathbb E_{N(d,\nu)}[\tanh y\,\operatorname{sech}^2y]\ge0$ for $d\ge0$.
Hence $|H_\nu'(d)|<1$ for every finite $d\ge0$: the upper bound is
strict since $\mathbb E\operatorname{sech}^2y<1$, and the pointwise
lower bound of $(1-m^2)(1-2m)$ is strictly greater than $-1$.
The strict derivative bound is uniform on the compact segment joining
$d$ and $\nu$.  Consistency gives $d-\nu=H_\nu(d)$ and therefore forces
$d=\nu$.

Finally $F(v)=\mathbb E_{N(v,v)}\tanh y$ satisfies $F(0)=0$ and
\[
 F'(v)=\mathbb E_{N(v,v)}[(1-m^2)(1-m)]\in[0,1)
 \qquad(v>0).
\]
The upper bound uses the same reflection identity.  Thus $v-F(v)$
is strictly increasing from zero to infinity, since $0\le F(v)\le1$.
It has a unique inverse at every $s>0$, with $s\le v_*(s)\le s+1$.
This proves continuity and the positive lower bound for $q$ on each
compact positive interval.  At $s=0$ there is no equality law with
$q>0$: its mean equation forces $d=0$, while
$\mathbb E_{N(0,q)}\tanh^2 y<q$.

We specify why the integral consequences extend as stated.  In the
smoothed integral the parameters are
$d=s+M$, $\nu=s+\sigma^2+q$, and the exponent is
$\psi=(a-\nu b)^2/[2\nu(\nu+q)]$.
With $\varphi_{d,\nu}$ the density of $N(d,\nu)$, the scalar
weight meant here is
\[
 G_{s,\sigma}(y)=
       \prod_{i=1}^n\varphi_{d,\nu}(y_i)\,\exp\{n\psi\}.
\]
When $\nu\ge\nu_{\min}>0$, Cauchy--Schwarz gives the uniform envelope
$G_{s,\sigma}(y)\le\exp(Cn-c\sum_i y_i^2)$. On a fixed empirical second-moment
ball, $M,q,a,d,\nu,\psi$ are continuous for weak convergence, since
only bounded or at most linearly growing functions are averaged.
The entropy variational formula and a finite covering by fixed
product-Gaussian neighborhoods give
\[
 \limsup_{n\to\infty}n^{-1}\log\int_{P_y\in\mathcal C}G_{s,\sigma}(y)\,\mathrm dy
 \le\sup_{P\in\mathcal C}\{\psi-D(P\Vert N(d,\nu))\}
\]
for each compact parameter set $\mathcal C$; Gaussian density ratios
on the second-moment ball tend uniformly to one at the exponential
rate when their parameters converge.

The rate is nonpositive, and
its zero set at $\sigma=0$ is exactly the curve just proved.
Upper semicontinuity therefore makes the rate strictly negative on
closed sets outside its prescribed neighborhood.  Compactness preserves
this negativity when $\sigma$ is sufficiently small.  The exterior of
the ball has any required negative rate by the envelope.
For the tail-square assertion insert
$\exp(\lambda\sum_i y_i^2\mathbf1_{|y_i|>R})$ into each of the finite
Gaussian comparisons, with fixed sufficiently small $\lambda>0$.

Its additional logarithmic moment tends uniformly to zero as
$R\to\infty$.  Removing the tilt on a positive tail-square level
produces strict negative rate.  These are precisely the three stated
integral conclusions, now with no subcritical assumption.
\end{proof}

\begin{proposition}[Positive-time field inputs]\label{prof:positive-stability}
Fix $0<s_1<T<\infty$.  There are positive constants
$\rho_0,\eta_0,\epsilon_A,c_0,c_*,a$, depending on this interval,
such that, at every deterministic $s\in[s_1,T]$, the probability under
the planted law, restricted to the original spectral event, of either
failure of \eqref{prof:eq26} or existence of $y$ with
\[
 \|F_{h_s}(y)\|\le\rho_0\sqrt n,\qquad q(y)<c_*
\]
is at most $e^{-an}$ for sufficiently large $n$.  The same assertion
holds uniformly on each fixed compact interval after decreasing the
residual threshold.

There are also absolute $c,C>0$ with the following small-time
refinement.  For each fixed $A_*>0$, at all sufficiently small
$r=\sqrt s>0$, failure of \eqref{prof:eq26} has probability at most
$e^{-A_*nr^3}$; its margins may depend on the scale.  For
$u\ge Cr$ in a sufficiently small fixed range, existence of an
approximate zero with $q(y)>u$ has probability at most $e^{-cnu^3}$.
The approximate-zero tolerance is fixed before taking $n\to\infty$
and may depend on $s,u,A_*$.  Uniformity is asserted only on compact
positive bands; finitely many or dyadic values of $u$ may be imposed
together.  In the small-time assertions a vanishing maximum diagonal
may be included in the external disorder restriction.
\end{proposition}

\begin{proof}
The scalar calculation separates the overlap from zero and places the
retained parameters near the positive-time equality curve. This gives
a norm--trace margin for nearby diagonals. We next check the conditional
finite-rank matrix correction and make that check simultaneous over
the diagonal neighborhood.

First retain the diagonal and gauge the planted spin to $\mathbf1$.
Then $J=W+\mathbf1\mathbf1^{\mathsf T}/n$ and
$h_s=s\mathbf1+\sqrt s\,g$. Until the final diagonal-removal
paragraph, use the temporary full-matrix residual and Hessian
\[
 \widetilde F_h(y)=y-h+(b(y)I-J)m(y),\qquad
 \widetilde{\mathcal H}(y,A)=I+A^{1/2}
 \left(b(y)I-J-\frac{2m(y)m(y)^T}{n}\right)A^{1/2}.
\]
Let \(E\subseteq\mathbb R^n\times\operatorname{Sym}_n\) be Borel
and independent of the observation \(h_s\). For \(\sigma>0\) set
\[
 I_{s,\sigma}(E)=(2\pi\sigma^2)^{-n/2}\int_{\mathbb R^n}
 \mathbf1_{\{(y,J)\in E\}}
 \exp\left[-\frac{\|\widetilde F_{h_s}(y)\|^2}{2\sigma^2}
                  +\frac{nb(y)^2}{2}\right]\,\mathrm dy.
\]
The exact Gaussian reduction, Lemma 5.3 of
\cite{AcceptedGap}, is valid for every positive $j$ and
here reads
\[
 \mathbb E I_{s,\sigma}(E)
 =\int\sqrt{\frac\nu{\nu+q}}\,G_{s,\sigma}(y)
 \mathbb P\bigl((y,J)\in E\mid
 Wm+\sqrt{s+\sigma^2}\,g'=y-d\mathbf1+bm\bigr)\,\mathrm dy,
\]
where $d=s+M$, $\nu=s+\sigma^2+q$ and $G_{s,\sigma}$ is the
explicit scalar weight in Lemma~\ref{prof:endpoint-scalar}.
Here \(g'\) is a fresh standard normal independent of \(W\).
The identity follows directly by inverting
$\nu I+mm^{\mathsf T}/n$; the determinant factor is
$\sqrt{\nu/(\nu+q)}$.

On $[s_1,T]$, discard a large second-moment ball's exterior, a fixed
neighborhood of the equality curve's complement, and the scalar
square-tail event.  Lemma~\ref{prof:endpoint-scalar} gives strictly
negative rates on these discarded parts and an arbitrarily small
positive rate for the total integral.  On the remaining parameters
$q\ge q_0>0$.  Choose $\eta_0$ and the diagonal Frobenius tolerance so
small that
\[
 \|A\|\chi_A\le(1+\eta_0)(1-q_0+\epsilon_A)
 \le1-q_0/2.
\]
The restricted fixed-parameter and bilinear estimates now apply.
We verify their conditional use and its positive margin.

Set $u=m/\sqrt{nq}$, $Z_1=(y-d\mathbf1)/\sqrt{n\nu}$ and
$P_u=I-uu^{\mathsf T}$.  Under the displayed conditioning, using a
fresh GOE $W_0$, put \(\ell_0=u^TW_0u\) and
\[
 \begin{split}
 P_uWP_u&=P_uW_0P_u,\\
 P_uWu&=\bar w+\tau P_uW_0u,\\
 \bar w&=\sqrt{q/\nu}\,Z_1-(a/\nu)u,
 \qquad \tau^2=(s+\sigma^2)/\nu,\\
 u^{\mathsf T}Wu&=\alpha_c
 =2(a+bq)/(\nu+q)
       +\sqrt{(s+\sigma^2)/(\nu+q)}\,\ell_0.
 \end{split}
\]
Writing out the three blocks gives the exact finite-rank difference
\[
\begin{split}
 W-W_0={}&u\bar w^T+\bar w u^T
 -(1-\tau)\{u(W_0u)^T+(W_0u)u^T\}\\
 &+\{\alpha_c+(1-2\tau)\ell_0\}uu^T.
\end{split}
\]
The scalar \(\ell_0\) has variance \(2/n\) and is smaller than
any fixed tolerance outside an exponentially small event.

Write \(S_A=S_A(1;W_0)\) and define
\(\mathcal Y(p)=S_A^{-1/2}A^{1/2}p\). Set
\[
 U=\mathcal Y(u),\quad T_1=\mathcal Y(W_0u),\quad
 \overline W=\mathcal Y(\bar w),\quad
 X_1=\mathcal Y(\mathbf1/\sqrt n).
\]
The exact whitened Hessian is
\[
\begin{split}
 S_A^{-1/2}\widetilde{\mathcal H}(y,A)S_A^{-1/2}
 ={}&I+(b-\chi_A)S_A^{-1/2}AS_A^{-1/2}
       -X_1X_1^T-2qUU^T\\
 &-\bigl[U\overline W^T+\overline WU^T
 -(1-\tau)(UT_1^T+T_1U^T)\\
 &\hspace{35mm}
       +\{\alpha_c+(1-2\tau)\ell_0\}UU^T\bigr].
\end{split}
\]
For each fixed \(y,A\), the vectors
\(\mathbf1/\sqrt n,u,Z_1\) are deterministic relative to \(W_0\).
The bilinear tests in \Cref{prof:endpoint-words} give
\[
\begin{split}
 \mathcal Y(p)^T\mathcal Y(r)&\simeq p^TAr,\\
 \mathcal Y(p)^TT_1&\simeq\chi_Ap^TAu,\\
 \|T_1\|^2&\simeq\chi_A+\chi_A^2u^TAu
\end{split}
\]
for this fixed finite vector list. In particular
\(\|X_1\|^2\simeq\chi_A\), rather than \(b\) for an arbitrary
\(A\). These are comparisons of actual Gram entries, with a
chosen fixed accuracy.

The comparison Gram for \(T_1\) is represented by
\(\chi_AU+\sqrt{\chi_A}N\), where \(N\) is an abstract unit
direction orthogonal to the other three images. This is a model
of Gram entries, not a coupling with a new random vector.
Set aside the small full-rank term \(b-\chi_A\) and the scalar
\(\ell_0\). Eliminating the \(N\) direction, whose Hessian pivot
is the identity, gives the ideal coefficient
\[
\begin{split}
 \beta_U
 &=2q-\frac{2a}{\nu}+\frac{2(a+bq)}{\nu+q}
       -2(1-\tau)\chi_A+(1-\tau)^2\chi_A\\
 &=q\left(2-\frac{\chi_A}{\nu}
       -\frac{2(a-\nu b)}{\nu(\nu+q)}\right)
\end{split}
\]
on \(UU^T\) in the form being subtracted. The final term in
the first line is the Schur correction.

At the equality-law limit, \(\sigma=0\), \(d=\nu=v=v_*(s)\),
\(a=vb\), and the limiting diagonal is \(A=V_y\), so
\(\chi_A=b\). Only in this limit is \(\beta_U=q(2-b/v)\).
Write
$b=\mathbb E\operatorname{sech}^2y$ and
$e=\mathbb E[\tanh^2 y\operatorname{sech}^2y]$ under $N(v,v)$.
Use the rescaled triple
\[
 M_1=\sqrt q\,\mathcal Y(u),\qquad
 G_1=\mathcal Y(Z_1)/\sqrt v.
\]
The ideal remaining matrix is identity minus the form on
$(X_1,M_1,G_1)$ with coefficient matrix \(C\) and limiting
equality-law Gram matrix \(H\):
\[
 C=\begin{pmatrix}1&0&0\\0&2-b/v&1\\0&1&0\end{pmatrix},\qquad
 H=\begin{pmatrix}
 b&e&-2e\\e&e&b-3e\\-2e&b-3e&b/v-2b+6e
 \end{pmatrix}.
\]
The Gram matrix here is the limit of the empirical \(A=V_y\)
comparison data, not the actual finite-sample Gram. Its entries
follow by Gaussian integration by parts; for example $\mathbb E[m(1-m^2)]=e$,
$v^{-1}\mathbb E[(y-v)m(1-m^2)]=b-3e$, and the double integration by
parts gives $b/v-2b+6e$.  Elementary column operations yield
\[
 \det(I-HC)=(q+2e)(q+e)^2>0.
\]
Subtracting the first vector alone is positive because its squared
length is $b<1$.

The remaining two-vector coefficient block has only
one positive eigenvalue.  Hence the final matrix can have at most one
nonpositive direction; its positive determinant rules that out.
The preceding orthogonal-direction Schur complement is therefore
positive as well.  Continuity on the compact equality curve gives a
fixed positive margin.

The fixed-parameter calculation is complete. To make it simultaneous,
we first quantify its continuity. The margin survives small errors in
all Gram entries, coefficients,
and $\chi_A-b$.  To check uniformity over nearby $A$, each tested
coordinate product is bounded by $C(1+y_i^2)/n$.  On $|y_i|\le R_1$,
its change under $A-V_y$ is at most
$C_{R_1}\|A-V_y\|_{\mathrm F}/\sqrt n$; the remaining part costs the
prescribed square-tail tolerance.  The $V_y$-weighted coordinate
products are bounded continuous functions, so their equality-law
limits follow from weak closeness.

This proves the fixed-$A$
conditional estimate with a uniform positive exponential rate.
The path margin \(S_A\succeq c_\kappa I\) returns the whitened
margin to \(\widetilde{\mathcal H}\), after the small
\((b-\chi_A)\) term is absorbed.
We pay for simultaneity only on the coordinates where a diagonal can
move appreciably. On the bounded conditional-matrix norm event choose
a fixed sup-norm mesh spacing $\delta$ so the $1/2$-H\"older error of
$\widetilde{\mathcal H}(y,A)$ is smaller than one quarter of the margin. Only
$\epsilon_A^2n/\delta^2$ sites can differ from $V_y$ by more than
$\delta$.  Mesh those sites and leave the others at $V_y$.
The logarithm of the number of approximants is at most
\[
 n\left[\mathfrak h(\alpha)+\alpha\log(C/\delta)+o(1)\right],
 \qquad \alpha=\epsilon_A^2/\delta^2,
\]
where $\mathfrak h(\alpha)\to0$ as $\alpha\downarrow0$.

Use twice the Frobenius radius in the fixed-parameter estimate and
then decrease $\epsilon_A$ to pay this mesh.  Thus the conditional
matrix failure, integrated against $G$, has strictly negative rate.
This explicitly extends the positive-time part of Lemma 5.5 of the
companion article; its subcritical small-$q$ argument is not used.

The estimates so far control a smoothed integral of bad configurations.
To exclude a bad approximate zero, a witness must force that integral
to be large enough. The deterministic local-volume estimate supplies
this implication, with determinant, smoothing, and residual tolerances
chosen in that order.

Let \(\mathcal D_{\eta_0}\) be the compact set of diagonal
\(0\le A\le(1+\eta_0)I\). Let \(\epsilon_*>0\) be the full
diagonal radius retained by the conditional calculation, and choose
\(c_E\) below half its matrix margin. Use the relaxed Borel set
\[
 E^{\rm mat}=\left\{(y,J):\exists A\in\mathcal D_{\eta_0},\
 \|A-V_y\|_{\rm F}\le\epsilon_*\sqrt n,\
 \lambda_{\min}(\widetilde{\mathcal H}(y,A))\le c_E\right\},
\]
intersected with the external spectral and norm restriction.
Compactness of the diagonal range and continuity make this set
Borel and independent of \(h_s\). The scalar exclusions and the
conditional matrix rate just proved give, for sufficiently small
\(\sigma>0\), uniformly on the fixed positive band,
\[
 \E I_{s,\sigma}(E^{\rm mat})\le e^{-a_0n}
\]
for some \(a_0>0\).

The stricter witness event uses
\[
 \|\widetilde F_{h_s}(y)\|\le\rho\sqrt n,\qquad
 \|A-V_y\|_{\rm F}\le\tfrac12\epsilon_*\sqrt n,\qquad
 \lambda_{\min}(\widetilde{\mathcal H}(y,A))\le\tfrac12c_E,
\]
with \(A\in\mathcal D_{\eta_0}\). It depends on the observation,
unlike \(E^{\rm mat}\). Keeping this same \(A\), for
\(d_y=\|y'-y\|/\sqrt n\) we have
\[
 \|V_{y'}-V_y\|_{\rm F}\le2\|y'-y\|,\qquad
 |b(y')-b(y)|\le2d_y,\qquad
 \left\|\frac{m(y')m(y')^T-m(y)m(y)^T}{n}\right\|\le2d_y.
\]
Thus the Hessian changes by at most \(6(1+\eta_0)d_y\).
Every point of the ball
\(\|y'-y\|\le r_{\rm loc}\sqrt n\) lies in the same
\(E^{\rm mat}\) fiber if
\[
 r_{\rm loc}\le
 \min\{\epsilon_*/4,\ c_E/[12(1+\eta_0)]\}.
\]

We next obtain the overlap lower bound needed for the determinant estimate.
A uniform lower bound on $q$ at any witness is available before
using the determinant estimate.  On the external norm event,
\[
 \|\tanh h_s\|\le\|m(y)\|+\|h_s-y\|
 \le C\sqrt{nq(y)}+\|\widetilde F_{h_s}(y)\|.
\]
Independence of the planted observation coordinates and a fixed
Bernoulli deviation bound give
$\|\tanh h_s\|\ge c\min(\sqrt s,1)\sqrt n$ outside a fixed
exponentially small event, uniformly on compact time intervals.
For instance, at small $s$, a fixed positive fraction of the
standard Gaussian coordinates lie in $[1/2,1]$, on which the
hyperbolic tangent has magnitude at least $c\sqrt s$.
On the fixed band let \(v_->0\) be a common observation lower
bound, and let the external matrix norm be at most \(K\).
Taking \(\rho\le v_-/2\) forces
\[
 q(y)\ge q_{\rm vol}:=
       \left(\frac{v_-}{2(K+2)}\right)^2>0.
\]

At such a witness the Jacobian is exactly
\[
 \widetilde F_{h_s}'(y)
 =I+(bI-W)V_y-\frac{\mathbf1\mathbf1^T}{n}V_y
                    -\frac{2mm^T}{n}V_y.
\]
The last two terms have rank at most two and norm at most three.
Since \(b\le1-q_{\rm vol}\), Lemma~\ref{prof:endpoint-logdet}
with \(\kappa=q_{\rm vol}\) gives its regularized determinant
upper bound with any needed fixed failure budget.
The deterministic local-volume lemma, Lemma 5.6 of
\cite{AcceptedGap}, has hypotheses $j,K<\infty$ and hence
applies at $j=1$ without modification.

With the regularized
Jacobian $Q=\widetilde F_{h_s}'(y)^{\mathsf T}
              \widetilde F_{h_s}'(y)+\gamma I$, its change of variable
$y'=y+\sigma Q^{-1/2}g$ proves a lower bound $e^{-\varepsilon n}$
for this bad integral whenever the residual is sufficiently small
and \(2\sigma/\sqrt\gamma\le r_{\rm loc}\).
Its losses per spin are bounded by
\[
 \frac{\varepsilon_1}{2}
 +C\left[\frac\rho\sigma+\frac\sigma\gamma+
       \left(\frac\rho\sigma\right)^2+
       \left(\frac\sigma\gamma\right)^2+
       \frac\sigma{\sqrt\gamma}\right]+o(1).
\]
Choose the target loss \(\varepsilon<a_0/2\), then the
determinant tolerance $\varepsilon_1$, then the regularizer $\gamma$,
then $\sigma$, and finally the residual tolerance $\rho$.

Markov's inequality now bounds the witness probability by
\(e^{-(a_0-\varepsilon)n}\), together with the assigned determinant
and observation exceptions. For low overlap, choose
\(q_{\rm rel}\) below the equality curve's minimum \(q\), use
the relaxed set \(E^{\rm low}=\{q\le q_{\rm rel}\}\) with the
same external restriction, and take
the witness threshold \(q\le q_{\rm rel}/2\).
Since \(|q(y')-q(y)|\le2d_y\), a radius
\(r_{\rm loc}\le q_{\rm rel}/4\) gives the same inclusion.
The scalar negative rate then replaces the matrix rate in the
same volume argument.

For the stronger small-time budgets use
\Cref{prof:scalar-diagnostics}, without the determinant weight.
Its scalar integral estimates permit any preliminary fixed
exclusion budget \(A_{\rm pre}nr^3\), restrict
$q\asymp_{A_{\rm pre}}r$, and give the outer rate $cnu^3$ for $u\ge Cr$.
Choose \(A_{\rm pre}>A_*\) to reserve rate for smoothing and
local volume. For an upper-tail witness \(q>u\), use the relaxed
set \(E_u^{\rm up}=\{q\ge u/2\}\) and radius
\(r_{\rm loc}\le u/4\). Taking \(u\ge2Lr\), with \(L\) from
the scalar proposition, gives a preliminary rate proportional to
\(n(u/2)^3\).

We compare the actual smoothed scalar weight with the root weight.
On a fixed band \(s\in[s_-,s_+]\), \(s_->0\), put
\(\nu_0=s+q\), \(\nu_\sigma=\nu_0+\sigma^2\), and
\(\mathfrak v=\langle(y-d)^2\rangle\). For fixed \(y\),
\[
 \frac1n\log\frac{G_{s,\sigma}(y)}{G_{s,0}(y)}
 =-\frac12\log\frac{\nu_\sigma}{\nu_0}
 +\frac{\sigma^2\mathfrak v}{2\nu_0\nu_\sigma}
 +\psi(\nu_\sigma)-\psi(\nu_0),
 \quad
 \psi(\nu)=\frac{(a-\nu b)^2}{2\nu(\nu+q)}.
\]
On \(\langle y^2\rangle\le R^2\), all numerators are bounded
and \(\nu\ge s_-\), so the absolute value is at most
\(C_{R,s_-,s_+}\sigma^2\). The uniform Gaussian envelope from
\Cref{prof:endpoint-scalar} makes the exterior integral have any
prescribed fixed negative rate when \(R\) is large.
Choose \(R\) first, then \(\sigma\) so this logarithmic loss is
smaller than one quarter of the reserved gap
\((A_{\rm pre}-A_*)r_-^3\), or than the reserved fixed multiple
of \(u_-^3\) for a finite outer-threshold family. Polynomial
factors are absorbed for large \(n\). This preserves every scalar
rate used here for the smoothed integral.

If a diagnostic itself uses \(\nu_\sigma\), give its thresholds
fixed slack first. On the same ball its scalar parameters change
by \(O_R(\sigma^2)\). On the retained scalar region,
\(q\ge c r_->0\), and the conditional column coefficients have
the same bound. The stricter smoothed bad set is then contained in the relaxed
unsmoothed one. This supplies the required set comparison as well
as the weight comparison.
For the matrix estimate the spectral indicator must remain attached
to the actual conditioned matrix. Denote that matrix by
\(J_\sigma=W_\sigma+\mathbf1\mathbf1^T/n\); the subscript records
the conditioning, not a replacement of the external disorder.
With \(u=m/\sqrt{nq}\), \(E=u^\perp\), and
\(e=P_E\mathbf1/\sqrt n\), its exact regression gives
\[
 J_\sigma|_E=W_E+ee^T,\qquad
 P_EJ_\sigma u
 =\frac q{\nu_\sigma}\frac{P_E(y-d\mathbf1)}{\sqrt{nq}}
      +\frac M{\sqrt q}e
      +\sqrt{\frac{s+\sigma^2}{\nu_\sigma}}\,w,
\]
where the GOE block \(W_E\) and \(w\sim N(0,I_E/n)\) are
independent. Its scalar block is the regression displayed earlier.
In decreasing eigenvalue order, compression and rank-one interlacing
give, on \(\mathcal G_{\rm spec}(J_\sigma)\),
\[
 \lambda_j(J_\sigma)\ge\lambda_j(W_E)
                         \ge\lambda_{j+2}(J_\sigma),
 \qquad 1\le j\le n-2.
\]
Thus every cumulative count changes by at most two, and the top
edge of \(W_E\) is bounded by that of this same full matrix.

Here is the precise spectral information used by the stability
calculation. Put \(\Delta_{\rm sp}=q^{8/5}\) and restrict
\(R=(2+O(q^2)-W_E)^{-1}\) to depths greater than
\(\Delta_{\rm sp}\). The inherited edge counts and bulk
semicircle limit give, uniformly for the \(O(q^2)\) shifts in
the stability Schur calculation,
\[
 \frac{\operatorname{Tr}R}{n}
       =1+O(\sqrt{\Delta_{\rm sp}})+o_n(1),\qquad
 \frac{\operatorname{Tr}R^j}{n}
       \asymp_j\Delta_{\rm sp}^{3/2-j}\quad(j\ge2).
\]
The first estimate follows by removing the edge contribution to
the first trace; dyadic counts give the upper estimates for higher
powers, and one band at depth comparable to \(\Delta_{\rm sp}\)
gives their lower estimates. The fixed-rank count error contributes
only \(O((n\Delta_{\rm sp}^j)^{-1})\). Likewise the restriction
to depths at most a fixed \(\zeta\) has first trace divided by
\(n\) at most \(O(\sqrt\zeta)+o_n(1)\).
In particular the first trace is normalized to one, as required
for the cancellation in \eqref{prof:eq11}. The other needed margin
uses the actual full spectrum:
\[
 (b+b^{-1})I-J_\sigma\succeq(1-o(1))q^2I,
 \qquad b+b^{-1}=2+q^2+O(q^3).
\]

Condition first on the spectrum of \(W_E\). On every spectrum
satisfying these inherited bounds, its Haar eigenframe and the
independent Gaussian column obey the tests in
\eqref{prof:eq11} at any prescribed fixed rate in \(nr^3\),
before imposing \(\mathcal G_{\rm spec}(J_\sigma)\).
The scalar diagnostics bound every regression coefficient uniformly
on the fixed positive band. Choosing their error tolerances first
and then \(\sigma\) small preserves the column tests and their
\(q^2\) Schur margin. Consequently the estimate is for
\[
 \int_{\mathcal S}G_{s,\sigma}(y)\,
   \mathsf Q_{y,\sigma}
      (\mathcal G_{\rm spec}(J_\sigma)\cap\mathcal F_y)\,dy,
\]
with the preliminary rate paying the scalar-weight and mesh losses.
Here \(\mathcal F_y\) is failure of the required Hessian margin.
No event for the fresh comparison matrix \(W_0\) replaces the
displayed indicator, and no division by its conditional probability
is used. This argument transfers the stability calculation, not
the sharper inverse-trace accuracy for \(K_o\), which is not
needed here. All smoothing choices are made on a fixed positive
band before taking \(n\) large.

We spell out the additional diagonal uniformity in this use.
At the given positive scale truncate \(A^{-1}-I\) above at the same level
\(q^{1-\eta}\) used for \(V_y^{-1}-I\), assigning that truncation
value at \(A_{ii}=0\). The truncated scalar function is Lipschitz on $[0,1+\eta_0]$; its
negative part is at most $\eta_0$.  Choosing $\eta_0$ small preserves
the norm bound for the truncated perturbation. Shrinking
$\|A-V_y\|_{\mathrm F}/\sqrt n$ makes its normalized first two
moments and all the fixed-axis quadratic and column tests as close
as required to those in \eqref{prof:eq11}.

For axes involving $y$,
first discard their square tails by the scalar diagnostics, then use
bounded site profiles on the remaining sites.  For $u=m/\sqrt{nq}$
one can directly use
$\|(B(A)-B(V_y))u\|\le C\epsilon_A/\sqrt q$.
All scales here are fixed before $n\to\infty$.  Thus the Schur
calculation retains its margin with the prescribed rate.
The original precision is larger. Its positive lower bound also
gives a positive symmetric-normalization margin: combine it with
the bound $A^{-1}-CI$, then take limits if some entries vanish.
The preceding sparse diagonal mesh costs an arbitrarily small
exponential rate and makes the estimate simultaneous.

The observation lower bound just used has a fixed rate independent
of sufficiently small $r$, so it pays any specified $A_*r^3$ after
shrinking the upper scale.  The determinant event can have an
arbitrarily large fixed rate by Lemma~\ref{prof:endpoint-logdet}.
Apply the local-volume argument in the same order, now taking its
loss below the available $A_*r^3$ or $cu^3$ rate.  This proves the
small-time assertions, including their order of tolerance choices.

Finally remove the diagonal.  On $\max_i|J_{ii}|\le\delta_D$ the
normalized residual changes by at most $\delta_D$, and the
symmetric Hessian by at most $(1+\eta_0)\delta_D$.
On fixed positive bands take a sufficiently small fixed $\delta_D$;
its failure has exponential probability.  In the small-scale
assertions keep the vanishing-diagonal event as an external
restriction and then, for each fixed band, take $n$ large enough.
For example, take \(\delta_D\le\rho/2\) and
\((1+\eta_0)\delta_D\le c_E/4\). Then a zero-diagonal
witness with residual at most \(\rho\sqrt n/2\), diagonal radius
\(\epsilon_*/2\), and Hessian threshold \(c_E/4\) implies the
full-matrix witness above.

For the compact path version write \(h_t=t\mathbf1+B_t\) in the
gauged planted process. On a mesh interval of length \(\Delta\),
the coordinate Brownian suprema have a fixed exponential moment
after division by \(\Delta\). Independence across coordinates gives
\[
 \Pr\!\left\{\sup_{t_j\le t\le t_{j+1}}
       \frac{\|B_t-B_{t_j}\|}{\sqrt n}>\kappa_0\right\}
 \le \exp\{-n(c\kappa_0^2/\Delta-C)\}.
\]
Choose \(\kappa_0,\Delta\) so this bound has the required fixed
rate and the Brownian increment plus the deterministic drift is
at most \(\rho/4\) in normalized norm. A finite union covers the
fixed compact interval. A path witness with residual
\(\rho\sqrt n/8\) then has residual below \(\rho\sqrt n/2\)
at the left mesh endpoint. The Hessian, overlap, and diagonal
predicate are independent of the field and persist there.
This proves the compact assertion with its reduced residual
threshold, also for a fixed finite family of scalar thresholds.
All previous exclusions estimated failure intersected with the
original spectral restriction, using interlacing rather than a
new typicality assertion for a conditional largest eigenvalue.
\end{proof}

\subsection{Residual identities and the complete derivative buffer}

\begin{proposition}[Endpoint posterior consequences]
\label{prof:endpoint-posterior}
Fix a matrix norm bound. For the covariance conclusion, assume
positive margins in \eqref{prof:eq26} and suppose that every point
in its small-residual region has \(q(y)\ge c_*>0\).
On the fixed-length ordinary and restricted inverse-word events
specified below,
\[
 \|\operatorname{Cov}_{\mu_h}(G,X)\|
 \le C\bigl(\operatorname{Var}_{\mu_h}(G)+\mathcal D_h(G)\bigr)^{1/2}
 \quad\text{for every }G.
\]
In particular $\|\Sigma_h\|\le C$.

For the following mean conclusion, only the whole-region root
stability implication
\[
 \|F_h(y)\|\le\rho_0\sqrt n
 \quad\Longrightarrow\quad
 \mathcal H(y,V_y)\succeq c_0I
\]
is required, together with ordinary word diagnostics. No lower
bound on \(q(y)\) is assumed in this mean conclusion.
For $N=\mu_h(f^2)>0$ and $\nu_h=f^2\mu_h/N$, and each sufficiently
small fixed $\zeta>0$,
\begin{equation}\label{prof:eq27}
 \frac{\|\mathbb E_{\nu_h}X-\mathbb E_{\mu_h}X\|}{\sqrt n}
 \le C\zeta+\frac{C_\zeta}{\sqrt n}
                   \left(1+\frac{\mathcal D_h(f)}N\right).
\end{equation}
The leading constant \(C\) depends on the fixed norm and stability
margins, but not on \(\zeta\).
Without the stability hypothesis, put
\[
 Q_\delta(h)=\sup_{\|F_h(y)\|\le\delta\sqrt n}q(y),
\]
with an empty supremum equal to zero. The same ordinary diagnostics give
\begin{equation}\label{prof:eq27-coarse}
 \frac{\|\mathbb E_{\nu_h}X-\mathbb E_{\mu_h}X\|}{\sqrt n}
 \le 2\sqrt{Q_\delta(h)}+
       \frac{C_\delta}{\sqrt n}
                   \left(1+\frac{\mathcal D_h(f)}N\right).
\end{equation}
For each fixed $\delta>0$ this
assertion holds for sufficiently large $n$ on the corresponding
ordinary diagnostics.  All events are independent of $f,G$ and the
unit directions used to test these inequalities.
\end{proposition}

\begin{proof}
The root and inverse statement, Lemma 5.7 of
\cite{AcceptedGap}, is deterministic for every finite positive
$j$ and applies here at $j=1$. Its root conclusion uses only the
whole-region implication at \(A=V_y\) stated above; the overlap
lower bound is used later only for the inverse recipe.
Denote the unique zero by $y_*$,
and set $m_* =\tanh y_*$, $q_*=n^{-1}\|m_*\|^2$.
It gives \(\|y-y_*\|\le C\|F_h(y)\|\) when
\(\|F_h(y)\|<\rho_0\sqrt n\), with
\(C=1+(\|J^\circ\|+3)/c_0\). The strict residual inequality
will be preserved by the inner buffer below.

We verify the precise diagnostic hypotheses of the companion residual
proof. Fix \(J,h\) throughout this deterministic subsection; all
spin expectations below are under \(\mu_h\), unless a displayed
square reweighting specifies otherwise. For a spin function define
\[
 d_i g(x)=\frac{g(x)-g(x^{(i)})}{2x_i},\qquad
 dg=\nabla g=(d_i g)_i.
\]
For a vector function, \((\nabla g)_{ji}=d_i g_j\).
These are discrete spin-flip derivatives, distinct from the
formal site derivatives in the recipes. Put
\[
 m^0=x,\quad b_0=0,\quad h^1=h+J^\circ x,\qquad
 m^l=\tanh h^l,\quad b_l=n^{-1}\sum_i(1-(m_i^l)^2),\quad
 R_l=m^{l-1}-m^l,
\]
\[
 h^{l+1}=h+J^\circ m^l-b_lm^{l-1}.
\]
For each fixed depth the primary differentiation estimates of
Lemma 6.1 of the companion article need only ordinary words, with
contraction coefficient one.  Explicitly their formal derivatives are
\[
 H_1=J^\circ,\quad M_0=I,\quad
 M_l=\operatorname{diag}(1-(m^l)^2)H_l,\quad
 H_{l+1}=J^\circ M_l-b_lM_{l-1}.
\]
The diagonal predictions of $H_l$ cancel by inclusion--exclusion over
originally adjacent pairs of each noncrossing pairing.

The word
estimates give bounded operator, diagonal $\ell^2$, and off-diagonal
entrywise $\ell^4$ norms.  Taylor's error in a coordinate difference
is at most $C|d_i h_j^l|^2$; summing its square uses that $\ell^4$
bound.  Differentiating $b_lm^{l-1}$ adds a rank-one term of bounded
Frobenius norm, since $\|db_l\|\le C_l/\sqrt n$.
Thus actual derivatives differ from $H_l,M_l$ by bounded Frobenius
errors.  Every single flip changes a fixed primary vector by at most
$C_l$ in Euclidean norm.  None of these finite recurrences uses
$j<1$ or convergence as $l\to\infty$.

We use the following finite recipe class, as in Definition~6.2 of
\cite{AcceptedGap}. Fix the primary depth and finitely many seeds
$s^1,\ldots,s^v\in\mathbb R^n$ with bounded Euclidean norms.
The seeds may depend on the fixed $(J,h)$ but are constant as functions
of the spin configuration $x$. Let $\theta(x)$ be a finite tuple of
scalar parameters with
$\sup_x(|\theta_\alpha(x)|+\|d\theta_\alpha(x)\|)\le C$
for every component; the averages $b_l$ may be among them.
Given auxiliary fields $y^1,\ldots,y^{a-1}$, a source has the form
\[
 w_i^a=\sum_{j=1}^v
   \phi_i^{a,j}((h_i^l)_{l\in I_a},\theta)s_i^j
 +\sum_{b<a}\psi_i^{a,b}((h_i^l)_{l\in I_a},\theta)y_i^b.
\]
All index sets are finite. The site coefficients and every fixed
derivative needed by the recipe are uniformly bounded on their argument
ranges and on the line segments joining evaluations. Deterministic
site labels are allowed with the same bounds. Define
\[
 y^a=J^\circ w^a
 -\sum_l\left(n^{-1}\sum_i\partial_lw_i^a\right)m^{l-1}
 -\sum_{b<a}\left(n^{-1}\sum_i\partial_{y^b}w_i^a\right)w^b.
\]
The partial $\partial_l$ differentiates only the explicit $h_i^l$
argument, keeping $\theta$ and all auxiliary arguments fixed;
$\partial_{y^b}w_i^a=\psi_i^{a,b}$, and an absent partial is zero.
A terminal vector has the same source form. When a source is reused
in a later source or terminal vector, substitute its source formula
without expanding the auxiliary fields.

A recipe is admissible from level $p$ when every explicit primary
argument in these source and terminal formulas has index at least $p$.
This restriction does not apply to scalar parameters or to the
correcting vectors $m^{l-1}$ in the auxiliary formula. All constants
may depend on the fixed recipe and its coefficient bounds. In particular,
the finite residual inductions below may enlarge the recipe before its
diagnostic length and the eventual dimension threshold are fixed.

Only the directional covariance argument uses an inverse recipe.
Fix $k\ge2$ and
$\Omega_{a\rho}=\{\max(\|R_k\|,\|R_{k-1}\|)\le a\rho\sqrt n\}$.
When several indices occur we write this set as
\(\Omega_{a\rho}^{(k)}\); the superscript is suppressed while
\(k\) is fixed.
The exact recurrence gives
\[
 F_h(h^k)=J^\circ R_k-b_{k-1}(R_{k-1}+R_k)
                  +(b_k-b_{k-1})m^k,
 \qquad |b_k-b_{k-1}|\le2\|R_k\|/\sqrt n.
\]
Choose \(\rho\) so the displayed recurrence puts
\(\Omega_{2\rho}\) strictly inside the root lemma's residual
region. Consequently on $\Omega_{2\rho}$,
\[
 \|h^k-y_*\|+\|m^{k-1}-m_*\|\le C\rho\sqrt n,
 \qquad |1-b_{k-1}-q_*|\le C\rho,
 \qquad q(h^k)\ge c_*.
\]
These constants are independent of $k$.  Define
\[
 \Phi(z;r,a)=\int_0^1e^{a(1-u^2)/2}
        \frac{\cosh(u(z-r))}{\cosh z\cosh r}\,\mathrm du,
 \qquad a=1-b_{k-1}-q_*,\qquad A_i=\Phi(h_i^k;(y_*)_i,a).
\]
Differentiation in $u$ proves the exact identity
\[
 [(z-r-a\tanh z)-a\partial_z]\Phi(z;r,a)=\tanh z-\tanh r.
\]
At $a=0$, $\Phi$ is the secant slope of $\tanh$.

Its positive integral
and the bounded derivatives of its logarithmic factors give
$0<\Phi\le e^{|a|/2}$ and bounded ratios of every fixed derivative to
$\Phi$.  The derivative recipes evaluate their coefficients at nearby spin
configurations as well as at the selected residual point. We therefore
establish the diagonal estimates on the whole outer buffer. Once the
finite recipe and its number of flips are fixed, the threshold in the
dimension will ensure that every evaluation stays in that buffer.

The preceding root-distance estimates imply, on the whole outer buffer,
\begin{equation}\label{prof:implicit-diagonal-buffer}
 \|A-V_{h^k}\|_{\mathrm F}\le C\rho\sqrt n,
 \quad |\chi_A-b_k|\le C\rho,
 \quad \|A\|\le e^{C\rho},
 \quad
 \|A\|\chi_A\le e^{C\rho}(1-c_*+C\rho)\le1-c_*/2.
\end{equation}
Choose $\rho$ small enough for the last inequality and for all stable
field tolerances.  This choice precedes the depth and recipe choices.
The inverse diagonal therefore belongs to $\mathcal A_\kappa$ with
$\kappa=c_*/2$ everywhere it will be evaluated.

For either $w_{0,i}=A_i e_i$, $\|e\|\le1$, or
$w_{0,i}=n^{-1/2}\partial_r\Phi(h_i^k;(y_*)_i,a)$, the implicit recipe is
\[
 w=w_0+A(y^{\mathrm{imp}}-c_1m_*),\qquad
 y^{\mathrm{imp}}=J^\circ w-\chi_A w-c_1m^{k-1},\qquad
 c_1=n^{-1}\sum_i\partial_k w_i.
\]
In this site partial, $y^{\mathrm{imp}},c_1$ and $a$ are held fixed.
After this one implicit solve, $y^{\mathrm{imp}}$ may be used as the
first auxiliary in an ordinary recipe, with $\sqrt n c_1$ as an
additional scalar parameter and $m_*/\sqrt n$ as an additional seed.
The corresponding coefficient and parameter bounds are required on
the flip buffers below; no such bounds are imposed on the arbitrary
exterior extension of the implicit solution.
Put $\ell_i=\partial_z\Phi/\Phi$, $u_0=w_0/A$ and
$p_0=A\partial_k u_0$.  Then $\partial_k w=\operatorname{diag}(\ell)w+p_0$
and elimination gives the actual finite-dimensional solve
\[
 \left[I-A(J^\circ-\chi_AI)
       +\frac1n A(m^{k-1}+m_*)\ell^{\mathsf T}\right]w
 =A\left[u_0-(m^{k-1}+m_*)n^{-1}\sum_i(p_0)_i\right].
\]
Put
\[
 \mathcal L=\chi_AI-J^\circ+
       (m^{k-1}+m_*)\ell^{\mathsf T}/n.
\]
Because $\|\ell+m_*\|\le C\rho\sqrt n$, the two-sided normalized
coefficient \(I+A^{1/2}\mathcal LA^{1/2}\), where the displayed solve is
\((I+A\mathcal L)w=\cdots\), differs in operator norm by \(O(\rho)\)
from \(\mathcal H(h^k,A)\). The normalized coefficient need not
be symmetric. The latter matrix has a fixed positive margin by
\eqref{prof:eq26}, so a perturbation smaller than half that margin
still has smallest singular value bounded below and a bounded
inverse.
The unnormalized inverse is bounded by
\[
 (I+A\mathcal L)^{-1}
 =I-A^{1/2}(I+A^{1/2}\mathcal LA^{1/2})^{-1}
             A^{1/2}\mathcal L.
\]
This proves
$\|w\|+\|y^{\mathrm{imp}}\|+\sqrt n|c_1|\le C$, uniformly in
$k$ and the unit seed.  Removing the negative rank term in
$\mathcal H(h^k,A)$ and replacing $b_k$ by $\chi_A$ also gives
\[
 S_A(1;J^\circ)\succeq cI.
\]
This separate adaptive stability condition is essential; the norm--trace
restriction alone was only used for fixed-parameter expectations.

Every coordinate-difference solve uses the same displayed coefficient
matrix evaluated at the flipped configuration.  All coefficient
differences multiply vectors at the common starting configuration.
Their squared column norms sum to a constant: for a coefficient $\psi$
and a fixed bounded-norm vector $v$,
\[
 \sum_{i,j}v_i^2|d_j\psi_i|^2
 \le\|v\|^2\max_i\sum_j|d_j\psi_i|^2\le C.
\]
The averaged coefficients have differences of order $n^{-1/2}$;
the factors $m^{k-1}+m_*$ have norm at most $2\sqrt n$.

Columnwise inversion therefore gives
$\|\nabla w\|_{\mathrm F}+\|\nabla y^{\mathrm{imp}}\|_{\mathrm F}
+\sqrt n\|dc_1\|\le C_k$, as in Lemma 6.3 of the companion article.
If the finished finite recipe uses at most $N_f$ successive flips,
the primary bounds put every required flip-neighborhood of
$\Omega_{3\rho/2}$ inside $\Omega_{2\rho}$ once
$n>(2N_f\max_l C_l/\rho)^2$. Thus all four bounds in
\eqref{prof:implicit-diagonal-buffer} and the adaptive inverse margin
hold at every derivative evaluation.  The recipe and its depth are
fixed before this threshold on $n$ is chosen.

We also fix the whole-cube convention for these local recipes.
Keep the actual implicit solution \(y^{\rm imp},c_1\) on
\(\Omega_{2\rho}^{(k)}\), assign arbitrary finite values to them
outside that set, and define \(w\) and all later ordinary
auxiliaries by their stated formulas. The implicit equation is
required only on the outer buffer. Every scalar test in a local
residual statement is supported on \(\Omega_\rho^{(k)}\).
Its discrete differences use only the fixed finite flip
neighborhoods already placed inside the buffers, so the exterior
values never enter a tested estimate.

We record the remaining closure of the residual proof to ensure that
no other inverse diagonal is being assumed.  Freezing parameters in
formal differentiation yields
\[
 \nabla w=A\nabla y^{\mathrm{imp}}+
       \operatorname{diag}(\partial_k w)H_k+E_1,
 \quad
 \nabla y^{\mathrm{imp}}=(J^\circ-\chi_AI)\nabla w-c_1M_{k-1}+E_2.
\]
Solving this equation introduces exactly the one factor
$K=(I-AJ^\circ+\chi_AA)^{-1}$ multiplying
$\operatorname{diag}(\partial_k w)H_k-c_1AM_{k-1}+E_1+AE_2$.
Later auxiliaries are ordinary, so error propagation multiplies by
ordinary words only.

The Taylor errors have bounded Frobenius norm
and bounded trace against the required words.  For the only
non-nuclear product error, the decisive bound is
\[
 \sum_{i\ne j}|P_{ji}|\,|d_js_i|\,Q_{ij}
 \le\|P_{\mathrm{off}}\|_{\ell^4}
       \|\nabla s\|_{\mathrm F}\|Q\|_{\ell^4}\le C,
 \qquad Q_{ij}=\sum_l|d_jh_i^l|.
\]
The diagonal part uses the primary diagonal $\ell^2$ estimate.
All other dropped terms are bounded in nuclear norm, using
$\|BC\|_{\mathrm{nuc}}\le\|B\|_{\mathrm F}\|C\|_{\mathrm F}$.

The formal trace cancellation in Lemma 6.4 of the companion article
is algebraic.  Its only cyclic node closes as
\[
 \mathfrak S_{\mathrm{imp}}
 =K_A(z;J^\circ)
       \{\operatorname{diag}(\partial_k w)\mathfrak H_{h^k}
                        -z^2c_1A\mathfrak S_{h^k}\}.
\]
Both terms, including the averaged correction, must be retained.
For each dependency path, inclusion--exclusion over its originally
adjacent contracted pairs cancels every complete noncrossing pairing.
There is no empty derivative path since seeds do not depend on $x$.

The contracted trace prediction vanishes as a formal series and
is a polynomial by Lemma~\ref{prof:endpoint-words}. Its evaluation
at \(z=1\) is therefore zero; this does not require convergence
of the inverse's formal series. Each resulting trace has either a bounded-norm small
vector as a diagonal factor or a scalar $O(n^{-1/2})$.
In the first case rotate the small diagonal to the beginning of
the trace. Cyclic invariance of the normalized prediction preserves
the canceled value, and Cauchy--Schwarz gives
\[
 \left|\operatorname{Tr}\bigl(\operatorname{diag}(u)
                    (Q-\mathsf P(Q))\bigr)\right|
 \le \|u\|_2\,
       \|\operatorname{diag}(Q-\mathsf P(Q))\|_2\le C.
\]
In the scalar case the diagonal word error has trace at most
\(C\sqrt n\), canceled by its \(O(n^{-1/2})\) coefficient.
These bounds remain uniform after a norm-one diagonal test.
Thus every small vector \(U\) in an ordinary recipe obeys the
following bound on the whole cube, and every local recipe vector
obeys it on \(\Omega_\rho^{(k)}\) with the stated flip buffers,
including the vector \(\partial_k w\):
\[
 \|U\|+\|\nabla U\|_{\mathrm F}
           +\|\operatorname{diag}\nabla U\|_1\le C.
\]
This proves the derivative assertion with only the single inverse
whose full buffer has already been checked.

We state the residual interfaces that now apply, specializing
Lemmas 6.5--6.6 of \cite{AcceptedGap} to coefficient one. Put
\[
 \mathcal E_h(G)=
       \bigl(\mu_h(G^2)+\mathcal D_h(G)\bigr)^{1/2},
 \qquad
 \mathcal D_h(G)=\mu_h\sum_i(1-(m_i^1)^2)(d_iG)^2.
\]
If \(U\) is a terminal vector of an ordinary recipe admissible
from level \(p\), then
\[
 |\mu_h[G R_p^TU]|\le C\mathcal E_h(G)
\]
for every scalar \(G\). For a terminal vector of a local recipe
admissible from level \(p\), with implicit starting index \(k\),
the same inequality holds when \(p\le k\) and \(G\) is supported
on \(\Omega_\rho^{(k)}\), with the buffers active.
For an ordinary recipe admissible from level \(p\) and a multiplier satisfying
\(\sup_x(|H_0(x)|+\|dH_0(x)\|)\le C_0\), the square-density
version is
\[
 |\mu_h[f^2H_0R_p^TU]|
 \le C\bigl(\mu_h(f^2)+\mathcal D_h(f)\bigr).
\]
Only ordinary recipes are used in this last version.

The induction can be checked directly. Its base identity is
\[
 \mu_h\sum_i(x_i-m_i^1)GU_i
      =\mu_h\sum_i(1-(m_i^1)^2)d_i(GU_i).
\]
The preceding diagonal $\ell^1$ bound and Cauchy--Schwarz control it
by \(C\mathcal E_h(G)\).
For the induction from $p$ to $p+1$ multiply the terminal vector by
$\Phi(h_i^p;h_i^{p+1},b_p-b_{p-1})$, add one ordinary auxiliary,
and use the displayed scalar identity for $\Phi$.

Here $|b_p-b_{p-1}|\le1$, so all required coefficient bounds hold.
The term $n(b_p-b_{p-1})\operatorname{av}(\partial_pU')$ cancels by
\[
 R_p^{\mathsf T}m^{p-1}
       =n(b_p-b_{p-1})-R_p^{\mathsf T}m^p.
\]
Only smaller residual indices remain; no further inverse is introduced.
The square-density version follows from
$d_i(f^2)=2f\,d_if-2x_i(d_if)^2$ and has right side
$C(\mu_h(f^2)+\mathcal D_h(f))$.  It uses ordinary recipes only.
All these operations are finite, so one fixed diagnostic length
covers the whole argument.

For the implicit recipe define the terminal residual
\[
 \mathcal T_k=
 \{h^k-h-(J^\circ-b_{k-1}I)m^k\}^Tw
       -(b_k-b_{k-1})nc_1.
\]
For every \(G\) supported on \(\Omega_\rho^{(k)}\),
\[
 |\mu_h[G\mathcal T_k]|\le C\mathcal E_h(G).
\]
Indeed, on the buffer the exact expansion is
\[
 \mathcal T_k=
 R_k^Ty^{\rm imp}+(\chi_A-b_{k-1})R_k^Tw
 -b_{k-1}R_{k-1}^Tw-c_1R_k^Tm^k.
\]
The local weak rule applies to the first three terms. In the last
one use target \(m^k/\sqrt n\) and multiplier \(\sqrt n\,c_1\),
whose size and discrete differences are bounded on the buffer.

We finish by spelling out selection and the uniformity in the tests.
Put $T_l=R_l^{\mathsf T}m^l/\sqrt n$.  The ordinary residual rule gives
$\mu_h(T_l^2)\le C_l$ and
\[
 \mathbb E_{\nu_h}(T_l/\sqrt n)^2
 \le\frac{C_l}{\sqrt n}\left(1+\frac{\mathcal D_h(f)}N\right).
\]
The exact telescoping identity
$\|R_l\|^2/n=b_l-b_{l-1}-2T_l/\sqrt n$
supplies a pair of consecutive small residuals. More precisely,
choose a fixed depth \(d_{\rm sel}\) with
\(\lfloor d_{\rm sel}/2\rfloor\rho^2/4>2\), and then a scalar
tolerance with \(2d_{\rm sel}\varepsilon<1\). If
\(\max_{l\le d_{\rm sel}}|T_l|/\sqrt n\le\varepsilon\), the
telescoping sum is less than \(2\), forcing a consecutive pair
whose two norms are at most \(\rho\sqrt n/2\).

Take a smooth \(\vartheta:[0,\infty)\to[0,1]\), equal to one
on \([0,1/4]\) and zero on \([1,\infty)\), and set
\[
 C_k=\vartheta\!\left(\frac{\|R_k\|^2}{\rho^2n}\right)
       \vartheta\!\left(\frac{\|R_{k-1}\|^2}{\rho^2n}\right),
 \qquad
 \widehat C_k=C_k\prod_{i=2}^{k-1}(1-C_i),
 \quad 2\le k\le d_{\rm sel}.
\]
Pointwise,
\[
 0\le\widehat C_k\le C_k\le1,\qquad
 \operatorname{supp}C_k\subseteq\Omega_\rho^{(k)},\qquad
 \|dC_k\|+\|d\widehat C_k\|\le C_\rho/\sqrt n.
\]
The derivative bounds follow from the primary differences and the
finite product rule. Also \(C_k=1\) when both residual norms are
at most \(\rho\sqrt n/2\), and
\[
 D:=1-\sum_{k=2}^{d_{\rm sel}}\widehat C_k
       =\prod_{k=2}^{d_{\rm sel}}(1-C_k)\in[0,1].
\]
The scalar moment bounds and a finite union bound give
\[
 \mu_h(D>0)\le C_\rho/n,
 \qquad
 \nu_h(D>0)\le\frac{C_\rho}{\sqrt n}
                      \left(1+\frac{\mathcal D_h(f)}N\right).
\]
These are the conclusions of the companion Lemma 7.1, with the
preceding derivation of its endpoint hypotheses.

Residual selection now turns the finite recipe estimates into posterior
bounds. The covariance argument uses its one controlled inverse; the
square-reweighted mean argument below uses only ordinary recipes.
This distinction also separates their diagnostic requirements.

We make the terminal implication explicit. On the selected buffer
put \(p_k=m^k-m_*\) and
\(L_*=\sqrt n(\chi_A-(1-q_*))\). Subtracting the root equation gives
\[
 h^k-h-(J^\circ-b_{k-1}I)m^k
 =h^k-y_*-a m^k-J^\circ p_k+(1-q_*)p_k.
\]
Using the implicit equations and the scalar identity for \(\Phi\)
then yields the exact identity
\[
\begin{split}
 \mathcal T_k={}&
 \sum_i\bigl[(h_i^k-(y_*)_i-a m_i^k)w_{0,i}
                         -a\,\partial_k w_{0,i}\bigr]\\
 &+(1-q_*-\chi_A)p_k^Tw-c_1R_k^Tp_k.
\end{split}
\]
The final term is controlled by the local weak rule with target
\(p_k/\sqrt n\) and multiplier \(\sqrt n\,c_1\). The remaining
scalar error satisfies pointwise
\[
 |(1-q_*-\chi_A)p_k^Tw|
 =\frac{|L_*|}{\sqrt n}|p_k^Tw|
 \le C_0\rho |L_*|.
\]
Here the root-distance and source-size bounds give \(C_0\)
depending only on the fixed norm and stability margins, not on
\(k\) or the selection depth. The earlier choice of \(\rho\)
includes \(C_0\rho<1/2\).

For the source \(w_{0,i}=n^{-1/2}\partial_r\Phi\), differentiation
of the scalar identity in \(r\), with \(a\) fixed, gives
\[
 [(z-r-a\tanh z)-a\partial_z]\partial_r\Phi
       =\Phi-\operatorname{sech}^2r.
\]
Thus the sum in the terminal identity is \(L_*\). Globally,
\(|L_*|\le C\sqrt n\) and \(\|dL_*\|\le C_k\).
Use the local test \(G=C_k^2L_*\), whose support lies in
\(\Omega_\rho^{(k)}\). The cutoff difference bound gives
\[
 \mathcal E_h(C_k^2L_*)\le C_{\rho,k}
                 +\|C_kL_*\|_{L^2(\mu_h)}.
\]
The tested scalar error is at most
\(C_0\rho\,\mu_h(C_k^2L_*^2)\). Absorbing it into the left side
of the terminal inequality proves
\[
 \|C_kL_*\|_{L^2(\mu_h)}\le C_{\rho,k}.
\]

For the source \(w_0=Ae\), \(\|e\|\le1\), the same sum is
\(p_k^Te\). Use the local test \(\widehat C_kG\).
Its support and difference bounds control its test norm by
\(C_\rho\mathcal E_h(G)\). For the scalar error use
\(\widehat C_k\le C_k\), the last \(L^2\) bound, and
Cauchy--Schwarz. The terminal and weak residual inequalities give
\[
 |\mu_h[\widehat C_kG(m^k-m_*)^Te]|
       \le C_{\rho,k}\mathcal E_h(G).
\]
Summing the ordinary weak rule with target \(e\) from level one
through \(k\) also controls the tested \((X-m^k)^Te\).

The deficit costs at most
$2\sqrt n\|G\|_2\mu_h(D>0)^{1/2}$.
This is the deterministic argument of Proposition 7.2, applied only
after checking its hypotheses on the selected regions and verifying
both of its actual source diagonals above.  It yields the asserted covariance inequality after
replacing $G$ by $G-\mu_h(G)$ and taking the supremum over $e$.
For $G=e^{\mathsf T}X$, $\mathcal D_h(G)\le1$; taking a top covariance
eigenvector gives $\|\Sigma_h\|\le C(\|\Sigma_h\|+1)^{1/2}$ and hence
a constant bound.

For the mean estimate, the ordinary residual rule suffices; no implicit
source is needed.  On each selected
region $m^k$ is within $C\zeta\sqrt n$ of $m_*$ if the residual
threshold was chosen below a fixed multiple of $\zeta$.
The ordinary square-density rule, summed with the constant target $e$
through $k$, bounds the tested difference between $X$ and $m^k$ by
$C_\zeta(1+\mathcal D_h(f)/N)$.  The deficit has the same bound.
This proves the conclusion of the companion Proposition 7.3 relative
to $m_*$.

Apply it also with $f=1$ and use the triangle inequality
to obtain \eqref{prof:eq27}.
If stability is absent, choose the iteration threshold so the exact
formula for $F_h(h^k)$ is at most $\delta\sqrt n$ on each selected
region.  Then $\|m^k\|\le\sqrt{nQ_\delta(h)}$, and the same argument,
once for $\nu_h$ and once for $\mu_h$, proves the coarser assertion.
This also covers an empty approximate-zero set, since then every
selected cutoff vanishes and the displayed deficit estimate controls
the entire measure.

All bounds use only norm bounds on seeds, so they hold for the whole
Euclidean unit ball simultaneously.  The finite recipe choices may
depend on the requested scalar tolerance, but the diagnostic event
quantifies over every bounded diagonal parameter after $J$ is observed.
It therefore does not depend on $f,G$, or the directional seed.
For covariance choose the stability margins and $c_*$, then $\rho$
and $\kappa$, then the selection depth and finite recipes, then their
word length, coefficient bounds and truncations, and finally the
size threshold guaranteeing every flip buffer.  For
\eqref{prof:eq27} and its coarser form, only the ordinary diagnostics
are required throughout.
\end{proof}

\subsection{Large observation time and the product endpoint}

\begin{proposition}[Terminal extension, uniformly in the initial spin law]
\label{prof:terminal-extension}
There are fixed $T,C,c>0$ and disorder events of ordinary probability
at least $1-e^{-cn}$ with the following property.  For every matrix in
these events, let $X$ have any probability law on the cube, possibly
depending on the matrix, and let the observation noises be independent
of $X$.  Set $h_T=TX+\sqrt T\,g$.  The reference Gibbs law $\mu_{h_T}$
has unscaled gap at least $1/2$ except with probability $e^{-cn}$.

On the original spectral event choose a fixed $c_+$ such that
$R=J+c_+I\succeq0$.  For one further time unit add observations with
drift $RX$ and covariance rate $R$, obtaining $h_{T+\lambda}$.
Define the reference laws
\[
 \mu_{h,\lambda}(x)\ \propto\
 \exp\left(\frac{1-\lambda}{2}x^{\mathsf T}J^\circ x
                         +h^{\mathsf T}x\right),
 \qquad 0\le\lambda\le1.
\]
The covariance and inverse unscaled gap of
$\mu_{h_{T+\lambda},\lambda}$ are bounded by $C$, simultaneously in
$\lambda$, outside an event of conditional probability $e^{-cn}$,
uniformly over the law generating $X$.  At $\lambda=1$ the reference
law is product.  These reference laws are the posteriors of the original
Gibbs prior evaluated at the observations; no assertion about the
posterior of the arbitrary law generating $X$ is made.
\end{proposition}

\begin{proof}
The signed curvature inequality, Lemma 8.1 of
\cite{AcceptedGap}, is deterministic for arbitrary symmetric
zero-diagonal interaction matrices with sufficiently small entries.
For interaction $B$, write
$U_{ij}=\tanh B_{ij}$, $Q_{ij}=U_{ij}^2$ (entrywise),
$r_4(U)=\max_i\sum_jU_{ij}^4$, and
$a_i=v_i d_i f$, where $v_i$ is the conditional spin variance.
It states
\[
 \mu_h[(-Lf)^2]\ge(1-Cr_4(U))\mathcal D_h(f)
       -\mu_h[a^{\mathsf T}Ua+C|a|^{\mathsf T}Q|a|].
\]
Its proof has no inverse-temperature parameter, so this is an exact
companion input.

Here is the endpoint matrix event, with exponential rather than only
asymptotic probability. Fix a curvature tolerance
\(\varepsilon>0\) so small that \((1+C)\varepsilon\le1/8\).
Choose a fixed global bound \(M\) for the row-square sums and
\(\|J^\circ\|\), with exponentially small exception. Next choose
an entry tolerance \(\delta>0\) small enough for the curvature
lemma and for \(M\delta/3\le\varepsilon/2\) and
\(CM\delta^2\le1/8\). Gaussian tails give
\(\max|J^\circ_{ij}|\le\delta\) outside probability
\(e^{-c_\delta n}\). For each fixed unit vector \(z\),
\[
 \mathbb P(|z^{\mathsf T}J^\circ z|>t)\le2e^{-nt^2/4},\qquad
 \mathbb P\left(\sum_{l\in S}(J^\circ_{il})^2>t\right)
 \le e^{-nt/4+|S|\log2/2}.
\]
A $1/4$-net of a $k$-dimensional unit sphere has size at most $9^k$.

Union these bounds over supports of size at most $\alpha n$ and the
$9^{\alpha n}$ nets.  Since
$n^{-1}\log\binom n{\lfloor\alpha n\rfloor}\to\mathfrak h(\alpha)$,
choosing $\alpha$ sufficiently small gives, with exponential probability,
\[
 \|(J^\circ)_{SS}\|\le\varepsilon/2,
 \qquad \max_i\sum_{l\in S}(J^\circ_{il})^2\le\varepsilon
 \quad(|S|\le\alpha n).
\]
These are the net and row arguments of Lemma 8.2 of the companion
article, now evaluated at variance $1/n$.

They work simultaneously for all $B=\theta J^\circ$, $0\le\theta\le1$.
Indeed $|\tanh z-z|\le|z|^3/3$ gives
$\|\tanh(\theta J^\circ)-\theta J^\circ\|\le M\delta/3$ by row sums;
$Q_{ij}\le(J^\circ_{ij})^2$ and $r_4(U)\le M\delta^2$.
By the preceding choice of \(\delta\), these errors fit the
\(\varepsilon\) and curvature budgets. We obtain uniformly in
$\theta$ bounded global norms of $B,U,Q$, small principal submatrix
norms of $U,Q$ on $|S|\le\alpha n$, and $Cr_4(U)\le1/8$.
The bounds hold pointwise over spin configurations and require no
union over configurations or initial laws.

The observation source law may differ from the Gibbs law whose gap is
being estimated. Throughout this argument the posterior is formed from
the Gibbs reference law. A norm estimate makes most coordinates strongly
biased for every source configuration; the small-submatrix bounds
control the remaining coordinates uniformly.

At time $T$ the observation field is $TX+\sqrt T\,g$.
During the extension it is
\[
 h_{T+\lambda}=TX+\sqrt T\,g+\lambda RX+R^{1/2}B_\lambda,
\]
where $B$ is a fresh standard Brownian motion independent of $X,g$.
Updating the original Gibbs prior by this likelihood subtracts
$\lambda x^{\mathsf T}Rx/2$ from its interaction; the diagonal is
constant on the cube.  This identifies its posterior with the stated
reference law.  The probability estimates below are uniform over the
possibly different law generating the observation source $X$.

On the norm event $\|R\|\le C$, the perturbation to $TX$ obeys
\[
 \sup_{0\le\lambda\le1}\|h_{T+\lambda}-TX\|
 \le C(1+\sqrt T)\sqrt n
\]
outside an event of probability $e^{-cn}$, independently of the law of
$X$.  The Gaussian norm tail controls $g$, and exponential Markov on
the sum of the coordinate Brownian suprema squared controls $B$.
Thus, for sufficiently large fixed $T$, fewer than $\alpha n/2$
coordinates have $|h_{T+\lambda,i}|<T/2$, uniformly in $\lambda$.
For every configuration $x$ simultaneously,
$\|(1-\lambda)J^\circ x\|\le M\sqrt n$, so at most
$16M^2n/T^2\le\alpha n/2$ coordinates have interaction field
larger than $T/4$.  Consequently the set where the full conditional
field has magnitude below $T/4$ has size at most $\alpha n$,
for every $x$ and $\lambda$.

Put $q_f(x)=\sum_i v_i(x)(d_i f(x))^2$ and
$\eta=\operatorname{sech}^2(T/4)$.  On the complement of that small
set, $\|a_{S^c}\|^2\le\eta q_f$, while $\|a\|^2\le q_f$.
The small-submatrix and global bounds therefore give pointwise
\[
 a^{\mathsf T}Ua+C|a|^{\mathsf T}Q|a|
 \le(1+C)\{\varepsilon+M(2\sqrt\eta+\eta)\}q_f.
\]
With \(\varepsilon\) already fixed, choose \(T\) large enough
both for the sparse-set bounds above and for
\((1+C)M(2\sqrt\eta+\eta)\le1/8\). The displayed expression
is then at most \(\mathcal D_h(f)/4\) after averaging.
The curvature inequality yields
$\mu_h[(-Lf)^2]\ge\mathcal D_h(f)/2$.
We finish by converting curvature to the gap.

By the finite-dimensional
spectral theorem each nonzero eigenvalue
of $-L$ is at least $1/2$.  The kernel consists of constants because
the Ising law is positive on every cube edge.  Thus
$\operatorname{Var}_{\mu_h}f\le2\mathcal D_h(f)$, uniformly in
$\lambda$ and $f$.  Applying this to unit linear tests also gives
$\|\Sigma_h\|\le2$.  This proves the proposition.
\end{proof}

We conclude by making explicit the probability interfaces used elsewhere.
For each fixed positive compact interval, combine
Propositions~\ref{prof:positive-stability} and
\ref{prof:endpoint-posterior} with the finite diagnostic events of
Lemma~\ref{prof:endpoint-words}.  Their ordinary exponentially small
exceptions remain exponentially small after planting: if
$L=Z(0)/\mathbb E Z(0)$, the critical overlap-binomial identity gives
$\mathbb E L^2\le C(n+1)$, for example by the pointwise Chernoff bound
$\mathbb P(\sum_i\varepsilon_i=k)\le2e^{-k^2/(2n)}$ and summing its
$n+1$ possible values.

Cauchy--Schwarz therefore pays only a polynomial
factor in transferring any disorder event.  It follows that, at each
deterministic positive time in that interval, $\|\Sigma_s\|$ is bounded
except with probability $e^{-cn}$ integrated under planting and the
original spectral restriction.  The terminal proposition gives the
corresponding terminal gap input.

For the macroscopic entropy argument, only the ordinary diagnostics in
\eqref{prof:eq27} and its $Q_\delta$ version are needed.  They are
field-independent and may be imposed externally on typical ordinary
disorder for each fixed tolerance.  Their exponent need not be compared
with the arbitrarily large multiplier $A_*$ in
Proposition~\ref{prof:positive-stability}.  On any compact positive
band with finitely many outer thresholds, use one residual tolerance
smaller than every threshold's permitted tolerance.

Existence of a
large-$q$ point at this tolerance implies existence at each
corresponding looser tolerance, so the outer rates remain valid.
Here is how the moving thresholds in the entropy argument fit this
finite construction. Let \(u_*>0\) be a fixed upper endpoint for
the small-range outer estimate, and on a fixed compact band put
\(R_b=C'A_*^{1/3}\sqrt b\), with \(A_*\ge1\) and \(C'\ge2C\).
Restrict the initial interval so that \(R_b\le u_*\).
For every moving threshold \(u=2^jR_b\le u_*\), choose a lower
point \(v\) in a fixed dyadic grid on that band with
\(u/2\le v\le u\). Then \(v\ge C\sqrt b\), and
\[
 \{Q_\delta(h_b)>u\}\subseteq\{Q_\delta(h_b)>v\},\qquad
 cv^3\ge(c/8)u^3.
\]
Only finitely many grid points occur on a fixed positive band.
For each point use the compact time subband on which
\(v\ge C\sqrt b\), and choose one residual tolerance for this
finite collection. Thresholds above \(u_*\) are grouped into the
single event \(Q_\delta>u_*\), using its fixed rate
\(e^{-cnu_*^3}\) and the bound \(Q_\delta\le1\).
Countably many bands and tolerances are handled by the diagonal
subsequence procedure in the entropy argument; no estimate is claimed
uniformly as positive time or a chosen residual tolerance tends to zero.
